\section{Conclusion}
\label{sec:conclusion}

This paper investigated a deceptively simple question: if LLM populations can develop social norms in the lab, why don't they in the wild? Our answer is that norm emergence is an \emph{environmental engineering problem}, not a model capability problem.

Through a $2\times2\times2\times2$ factorial experiment ($N=480$), we established three empirical findings. First, communication structure is the dominant environmental determinant ($\eta^2 = 0.31$, $p < 0.001$), exceeding feedback ($\eta^2 = 0.14$), stability ($\eta^2 = 0.09$), and task framing ($\eta^2 = 0.06$). Second, feedback loops are necessary but not sufficient: they amplify convergence when communication infrastructure exists but cannot drive emergence alone. Third, our wild-type conditions quantitatively reproduce the Tsinghua 93\% non-response baseline, confirming that the lab-wild divergence is an environmental artifact rather than a model property.

We introduced NEEM (the Norm Emergence Environmental Model) as a predictive framework and translated it into three actionable platform design principles: shared context infrastructure, guaranteed message delivery with bounded latency, and population stability mechanisms. These interventions address the most common multi-agent failure mode (coordination failure, 36.9\% per MAST \citep{deepmind2025mast}) at lower cost than model capability upgrades.

For the CHI community, our findings reframe a central design challenge. As multi-agent LLM systems become a standard interaction paradigm, the question is no longer "can we build models that develop social norms?" but "what infrastructure should we build to make prosocial norm emergence likely?" Communication infrastructure—where agents can reliably see each other, remember shared history, and interact with bounded latency—is the primary lever. This is a software engineering and HCI design problem as much as an AI capability problem.

Future work should explore human-AI mixed norm formation (RQ4), the minimal identity framework for coordination (RQ5), and field deployment of NEEM-informed platform design. All experimental materials, code, and data will be released upon publication.
